\documentclass[12pt]{article}

\usepackage[a4paper,margin=1in]{geometry}
\usepackage{setspace}
\usepackage{titlesec}

\setstretch{1.5}
\titleformat{\section}{\large\bfseries}{}{0pt}{}

\begin{document}

\begin{center}
{\Large \textbf{Report on Activities and Functions of the Department of Statistics}}\\[0.2cm]
{\large Government of Kerala and Government of India}

\vspace{0.5cm}

\textbf{Submitted by}\\
\textbf{Aishwarya J S}
\end{center}

\section*{Introduction}

Statistics is essential for planning, governance, policy formulation, and socio-economic development. It enables governments to collect, analyze, interpret, and disseminate reliable information for evidence-based decision-making. In Kerala, the Department of Economics and Statistics (DES) performs this role at the state level, while the Ministry of Statistics and Programme Implementation (MoSPI) is responsible for official statistics at the national level. Together, they provide timely and accurate data for budgeting, development planning, monitoring public programmes, and evaluating welfare schemes.

\section*{Government of Kerala -- Department of Economics and Statistics}

The Department of Economics and Statistics functions under the Planning and Economic Affairs Department of the Government of Kerala. It is the state's nodal agency for collecting, compiling, analyzing, interpreting, and publishing statistics related to agriculture, industry, population, prices, and other sectors of Kerala's economy. Its headquarters is located at Vikas Bhavan, Thiruvananthapuram.

The department has a structured administrative network consisting of the Directorate, District Statistical Offices headed by Deputy Directors, and Taluk Statistical Offices headed by Taluk Statistical Officers. This system ensures efficient collection and dissemination of statistical information across Kerala.

DES originated from an Agricultural Statistics Survey conducted by Travancore University in 1949. It later became the Board of Statistics in 1951, expanded after Kerala's formation in 1956, became the Bureau of Economic Studies in 1958, and was renamed the Directorate of Economics and Statistics in 1980.

Its major activities include agricultural and crop estimation surveys, land use surveys, agricultural input surveys, the Agriculture Census every five years, preparation of district statistical handbooks, statistical abstracts, socio-economic reports, and price indices.

At present, DES prepares Gross State Domestic Product (GSDP), Net State Domestic Product (NSDP), and Per Capita Income estimates. It publishes the annual Economic Review, monitors Sustainable Development Goal (SDG) indicators, collects Consumer Price Index (CPI) data, and conducts digital socio-economic surveys for decentralized planning.

\section*{Government of India -- Ministry of Statistics and Programme Implementation (MoSPI)}

MoSPI is the central statistical authority of the Government of India. Established on 15 October 1999, it was formed through the merger of the Department of Statistics and the Department of Programme Implementation.

The National Statistical Office (NSO), under MoSPI, coordinates India's official statistical system by developing statistical standards, methodologies, classifications, and data dissemination practices while coordinating statistical activities among central ministries and states.

MoSPI conducts nationwide surveys on employment, education, health, consumption, and social conditions. It prepares GDP estimates, National Income Accounts, CPI, Index of Industrial Production (IIP), and other major economic indicators.

The ministry conducts the Periodic Labour Force Survey (PLFS), Annual Survey of Industries (ASI), and household surveys. It releases monthly CPI and IIP reports, quarterly and annual GDP estimates, monitors Central Sector Infrastructure Projects, and evaluates government programmes. Its eSankhyiki platform provides public access to official datasets and statistical publications.

\section*{Interlinkage between Kerala DES and MoSPI}

Kerala DES and MoSPI work together within India's two-tier statistical system. Kerala DES collects primary state-level data, while MoSPI provides standardized survey methods, definitions, and classifications. The data collected in Kerala contributes to national statistical estimates and supports evidence-based governance at both state and national levels.

\section*{Conclusion}

The Department of Economics and Statistics and MoSPI play a significant role in India's development planning and administration. DES focuses on state-level surveys, economic analysis, price monitoring, and planning, while MoSPI coordinates the national statistical system, publishes official economic indicators, conducts nationwide surveys, and monitors government programmes. Together, they ensure reliable statistics for policymaking, research, governance, and sustainable socio-economic development.

\end{document}
